\section{Khối vi điều khiển Arduino Pro Mini}
Arduino Pro Mini là trung tâm xử lý của node. Trên sơ đồ, Arduino Pro Mini được đặt ở khối vi điều khiển và kết nối ra các header để thuận tiện đấu dây, nạp chương trình và kiểm thử. Bo mạch này có kích thước nhỏ, dễ lập trình trong môi trường Arduino IDE hoặc PlatformIO, phù hợp với các ứng dụng cảm biến đơn giản.

\subsection{Vai trò của vi điều khiển}
Arduino Pro Mini thực hiện các nhiệm vụ sau:
\begin{itemize}
    \item Khởi tạo cảm biến và ngoại vi.
    \item Đọc giá trị cảm biến theo chu kỳ.
    \item Lọc nhiễu và quy đổi dữ liệu sang đơn vị sử dụng.
    \item Đóng gói dữ liệu để gửi qua LoRa.
    \item Nhận lệnh điều khiển từ gateway.
    \item Điều khiển relay và bộ phát hồng ngoại.
    \item Gửi phản hồi trạng thái sau khi thực hiện lệnh.
\end{itemize}

Do tài nguyên hạn chế, chương trình trên Arduino cần được tổ chức theo vòng lặp trạng thái rõ ràng. Các thao tác đọc cảm biến, truyền LoRa và điều khiển thiết bị không nên chặn hệ thống quá lâu. Nếu sử dụng delay quá nhiều, node có thể phản hồi lệnh chậm hoặc bỏ lỡ gói tin từ gateway.

\subsection{Bố trí header}
Sơ đồ khối vi điều khiển có các header P1, P2 và P4 đưa các chân quan trọng ra ngoài. Việc sử dụng header giúp mạch linh hoạt hơn trong quá trình phát triển. Khi cần thay đổi loại cảm biến hoặc kiểm tra tín hiệu, người dùng có thể đo trực tiếp trên header mà không phải can thiệp vào mạch in.

\begin{longtable}{|p{3cm}|p{4.2cm}|p{5.5cm}|}
\hline
\textbf{Header} & \textbf{Nhóm chân} & \textbf{Chức năng} \\
\hline
P1 & \texttt{RAW}, \texttt{GND}, \texttt{RST}, \texttt{VCC}, \texttt{DIO0}, \texttt{SCK}, \texttt{MISO}, \texttt{MOSI}, \texttt{NSS} & Cấp nguồn, reset và giao tiếp SPI với LoRa. \\
\hline
P2 & \texttt{UITXo}, \texttt{UIRXo}, \texttt{GND}, \texttt{Po}, \texttt{A}, \texttt{TX}, \texttt{RL1}, \texttt{RL2} & Giao tiếp UART, tín hiệu cảm biến analog, IR và relay. \\
\hline
P4 & \texttt{SCL}, \texttt{SDA}, \texttt{GND}, \texttt{A6}, \texttt{A7} & Kết nối I2C và analog mở rộng. \\
\hline
\end{longtable}

\subsection{Quản lý chân vào ra}
Việc quản lý chân cần được cố định trong firmware bằng các hằng số rõ ràng. Ví dụ, chân \texttt{RL1} điều khiển relay 1, \texttt{RL2} điều khiển relay 2, \texttt{TX} điều khiển phát IR, \texttt{Po} đọc pH và \texttt{A} đọc cảm biến ánh sáng. Với cách đặt tên này, khi đọc mã nguồn có thể đối chiếu trực tiếp với sơ đồ nguyên lý.

\begin{lstlisting}[language=C,caption={Ví dụ định nghĩa chân trên node Arduino}]
#define PIN_RELAY_1   7
#define PIN_RELAY_2   8
#define PIN_IR_TX     5
#define PIN_PH        A1
#define PIN_LIGHT     A0
#define PIN_TDS       A2
#define PIN_LORA_NSS  10
#define PIN_LORA_RST  9
#define PIN_LORA_DIO0 2
\end{lstlisting}

Trong hệ thống thực tế, các chân phải được kiểm tra lại theo mạch in cuối cùng. Nếu sơ đồ nguyên lý, layout PCB và firmware không thống nhất, hệ thống có thể hoạt động sai dù từng khối riêng lẻ không lỗi. Vì vậy, bảng ánh xạ chân là một tài liệu bắt buộc đi kèm phần cứng.

\subsection{Chu kỳ hoạt động của node}
Chu kỳ hoạt động cơ bản của node gồm bốn pha. Pha thứ nhất là đọc cảm biến. Pha thứ hai là xử lý và đóng gói dữ liệu. Pha thứ ba là gửi dữ liệu qua LoRa. Pha thứ tư là lắng nghe lệnh điều khiển từ gateway trong một khoảng thời gian ngắn. Chu kỳ này lặp lại liên tục.

\begin{figure}[h!]
    \centering
    \begin{tikzpicture}[node distance=1.3cm, every node/.style={font=\small}]
        \node[draw, rounded corners, minimum width=3.2cm, minimum height=0.9cm] (s1) {Đọc cảm biến};
        \node[draw, rounded corners, right=of s1, minimum width=3.2cm, minimum height=0.9cm] (s2) {Lọc và quy đổi};
        \node[draw, rounded corners, right=of s2, minimum width=3.2cm, minimum height=0.9cm] (s3) {Gửi LoRa};
        \node[draw, rounded corners, below=of s2, minimum width=3.2cm, minimum height=0.9cm] (s4) {Nhận lệnh};
        \draw[->, thick] (s1) -- (s2);
        \draw[->, thick] (s2) -- (s3);
        \draw[->, thick] (s3) |- (s4);
        \draw[->, thick] (s4) -| (s1);
    \end{tikzpicture}
    \caption{Chu kỳ hoạt động cơ bản của node}
    \label{fig:node_cycle}
\end{figure}

\subsection{Quy tắc an toàn khi khởi động}
Khi node vừa cấp nguồn, trạng thái relay cần được đặt về mức an toàn. Nếu relay điều khiển quạt hoặc đèn, hệ thống nên mặc định tắt thiết bị cho đến khi nhận được cấu hình hoặc lệnh hợp lệ. Với IR điều hòa, node không phát lệnh khi khởi động để tránh thay đổi trạng thái điều hòa ngoài ý muốn.

Trình tự khởi động đề xuất gồm: cấu hình chân output, đặt relay về OFF, khởi tạo Serial để debug, khởi tạo cảm biến, khởi tạo LoRa, gửi gói trạng thái đầu tiên và bắt đầu vòng lặp chính. Nếu LoRa không khởi tạo thành công, node vẫn có thể đọc cảm biến và báo lỗi qua Serial hoặc LED trạng thái.
